\chapter{Dei ti namna}

\begin{motto}
Eit fag med eit fundament siterer resultat.\\ Eit fag utan eitt siterer namn.
\end{motto}

Vi har til no rive ned tesane éin for éin: den tomme formelen, den vitskaplege ambisjonen, det vrange problemet, det refleksive atelieret, den semantiske vendinga. Men ein innvending står att, og ho er den mest nærliggjande lesaren kan reise. Om faget verkeleg er så grunnlaust som boka hevdar, korleis kan det då ha frambore tenkjarar av eit slikt format? Korleis kan eit fag utan fundament ha produsert namn som er siterte titusenvis av gonger, lesne på kvart einaste designstudium i verda, omsette til tjue språk? Eit fag utan grunn, lyder innvendinga, ville ikkje ha hatt nokon å sitere. At det har det — at det har sine klassikarar, sine grunnleggjande verk, sin kanon — må då vere prov på at grunnen er der likevel, berre at boka nektar å sjå han.

Dette kapitlet tek innvendinga på det største alvor, og snur henne. Vi skal gå gjennom dei mest siterte figurane i heile designforskinga — målt så ærleg som offentlege siteringstal tillèt — og spørje for kvar einaste ein: kva hevda han, kvifor vart det sitert, og er det eit fundament? Påstanden boka skal forsvare er hard, men presis. At desse namna er store, er sant. At dei er glimrande, er for dei fleste sant. Men det dei vart siterte \emph{for}, var i nesten kvart tilfelle ei innflytelsesrik \emph{meining} om kva design er — ikkje ei prøvd, kumulativ kunnskap om korleis form verkar. Og når toppen av eit fag er bygd av meiningar, er ikkje toppen eit prov på fundamentet. Han er det mest presise målet på fråvêret av det.

\section{Kva siteringstala faktisk seier}

Lat oss byrje med tala, for dei er meir avslørande enn både kritikarane og forsvararane av faget plar innrømme. Slår ein opp dei mest siterte forfattarane som vert rekna til designfeltet, teiknar det seg eit mønster så reint at det nesten er pinleg.

Den mest siterte er Donald Norman, med over hundre og sytti tusen siteringar og ein h-indeks kring hundre og seksten — tal som plasserer han mellom dei høgst siterte levande forskarane i kva fag som helst. Klaus Krippendorff følgjer like etter, med over hundre og tjue tusen. Herbert Simon, om ein reknar heile forfattarskapen, ligg endå høgare. Så fell det bratt: Nigel Cross kring førti tusen, Kees Dorst kring tjue tusen, Victor Margolin kring åtte tusen, Bryan Lawson, Harold Nelson og Erik Stolterman, Ezio Manzini lengre nede — fleire av dei kjende, kanoniserte verka med eit lågare firesifra tal kvar.

Sjå no nøye på dei tre på toppen, for dei avslører heile saka. Norman sine hundre og sytti tusen siteringar er ikkje siteringar av design. Brorparten av dei kjem frå kognisjonsvitskapen og menneske–maskin-interaksjonen: arbeid om merksemd, om minne, om mental representasjon, frå den delen av karrieren hans som låg \emph{før} han vart designtenkjar, og som vart sitert fordi det var prøvbar psykologi. Den same historia gjeld Krippendorff. Av hans hundre og tjue tusen siteringar kjem talet store fleirtalet frå éi bok — \emph{Content Analysis}, ein metodologi for innhaldsanalyse i samfunnsvitskapen, med over tretti tusen siteringar åleine. Den boka er sitert så mykje fordi ho er ein \emph{metode som verkar}: ho gjev deg ein prosedyre, ein reliabilitetskoeffisient, ein måte å rekne om to forskarar kodar likt. Og Simon — den tredje — vart Nobelprisvinnar i økonomi og Turing-prisvinnar i informatikk, og dei aller fleste av siteringane hans gjeld avgrensa rasjonalitet og kunstig intelligens, ikkje design.

Her er funnet, og det er øydeleggjande for innvendinga. Dei tre mest siterte «designtenkjarane» er mest siterte for arbeid som ikkje er design, og som — det er poenget — er prøvbart. Norman vart stor på eksperimentell psykologi. Krippendorff vart stor på ein etterprøvbar metode. Simon vart stor på formaliserbare modellar av avgjersler. I det augneblinket kvar av dei vende seg mot design og skreiv det verket designfaget kanoniserer dei for — \emph{The Design of Everyday Things}, \emph{The Semantic Turn}, \emph{The Sciences of the Artificial} sitt designkapittel — datt etterprøvbarheita ut, og att stod ei velformulert meining. Dei bar med seg autoriteten frå dei prøvbare faga inn i det uprøvbare, og faget tok autoriteten som om han følgde med på lasset. Det gjorde han ikkje. Ein h-indeks tent på psykologi garanterer ingenting om kvaliteten på ein designfilosofi. Faget siterer namnet og trur det siterer grunnen.

\section{Dei som hevda at design \emph{ikkje} kan grunnleggjast}

Gå så til den indre kjernen — dei som er store \emph{innanfor} design, ikkje på lån utanfrå. Den fyrste gruppa deler ei merkeleg felles læresetning: dei vart mest innflytelsesrike for å ha argumentert, kvar på sitt vis, for at design er ein kunnskapsform som ikkje lèt seg redusere til den vitskaplege.

Nigel Cross er den fremste. Hans «designerly ways of knowing» — fyrst ein artikkel i 1982, seinare ei bok som åleine er sitert kring åtte tusen gonger — er kanskje den mest siterte enkeltformuleringa i heile feltet \autocite{cross1982,cross2006}. Gjev han den sterkaste forma fyrst, for han fortener det. Cross hevda at designarar har ein eigen, legitim måte å kjenne verda på, skild frå humanioras og naturvitskapens: dei tenkjer gjennom skisser og prototypar, dei går fram ved å føreslå løysingar heller enn å analysere problem, dei les form direkte. Som \emph{observasjon} er dette skarpt og truleg sant; designarar gjer verkeleg dette. Men spør kva slag påstand «designerly ways of knowing» er, og vanskane kjem. Er det ein empirisk påstand om korleis designarar faktisk tenkjer — då måtte han kunne prøvast, og vere feilbarleg, og Cross sjølv samla protokollstudiar som peikte i mange retningar. Eller er det ein \emph{normativ} påstand om at denne måten å kjenne på er like gyldig som dei andre — då er det ikkje forsking, men eit forsvarsskrift. Cross gleid mellom dei to, og det er nett glidinga som gjorde formuleringa så brukbar. Ho kunne lesast som vitskap når faget ville verke vitskapleg, og som frigjering frå vitskapen når faget ville sleppe krava hennar. Eit fundament kan ikkje vere begge delar. «Designerly ways of knowing» vart sitert åtte tusen gonger ikkje fordi det var prøvd, men fordi det gav kvar lesar nett det han trong: ei lærd formulering av kjensla av at design er noko eige. Det er innflytelse, og det er ikkje kunnskap.

Kees Dorst arvar og raffinerer dette. Saman med Cross gav han faget «co-evolution» av problem og løysing — tanken om at designaren ikkje løyser eit fast problem, men lèt problem og løysing utvikle seg saman til dei passar \autocite{cross2001discipline}. Seinare gav han det «frame creation» og «the core of design thinking» \autocite{dorst2011core,dorst2015frame}. Den sterkaste forma: dette er presise skildringar av noko verkeleg i designprosessen, og dei resonnerer fordi kvar praktiserande designar kjenner dei att. Men kjenne att er ikkje det same som å prøve. Co-evolusjon er ein \emph{metafor} lånt frå biologien, og som vi alt har sett, lånt utan det biologien byggjer han på — utan eit tilpassingslandskap, utan ein måte å måle «fit», utan noko som gjer at to forskarar kunne usemje seg presist om kor godt eit problem og ei løysing har ko-evolvert. Det er eit bilete på prosessen, ikkje ein teori om han. Og bilete kan ikkje falsifiserast; dei kan berre verke meir eller mindre treffande. Dorst sine kring tjue tusen siteringar er siteringar av treffande bilete.

Bryan Lawson høyrer til same familien. \emph{How Designers Think}, sitert kring to tusen gonger, samla tiår med observasjon av designarar i arbeid \autocite{lawson1980}. Boka er klok, rik på døme, og fullstendig ærleg om det ho er: ei skildring. Lawson hevda aldri å gje faget ein teori med prediktiv kraft; han skildra korleis designarar synest å gå fram. Det er kanon-statusen som er avslørande, ikkje boka. Eit fag som gjer ei \emph{skildring} av praksis til eit av sine grunnleggjande verk, fortel deg at det ikkje har noko meir grunnleggjande å setje der.

Sjå mønsteret i heile denne gruppa. Cross, Dorst, Lawson — dei mest siterte stemmene om sjølve designtenkinga — vart store for å ha skildra og forsvart ein kunnskapsform, ikkje for å ha bygd henne ut til noko prøvbart. Det dei deler, er ein vakker observasjon opphøgd til identitet. Og ein identitet kan ein sitere; ein kan ikkje byggje vidare på han, fordi det ikkje finst noko å korrigere.

\section{Refleksjonen, det vrange og meininga}

Tre av dei aller mest siterte verka i faget har vi alt felt dom over kvar for seg, men dei må møtast att her, fordi dei saman utgjer sjølve fundamentet faget plar peike på når det vert utfordra. Det er talande at når ein bed ein designforskar om feltet sine grunnsteinar, kjem nett desse tre.

Donald Schön sin «reflective practitioner» er den eine \autocite{schon1983,schon1987educating}. Den sterkaste forma er sterk: mot ein teknisk rasjonalitet som ville pressa all profesjonell kunnskap inn i forma «bruk regel på tilfelle», viste Schön at dyktige utøvarar tenkjer \emph{i} handlinga, fører ein samtale med situasjonen, oppdagar problemet medan dei løyser det. Dette er ein verkeleg innsikt, og han har frigjort generasjonar av praktikarar frå ein falsk modell av seg sjølve. Men sjå kva slag innsikt det er. «Refleksjon-i-handling» skildrar at noko skjer; det forklarer ikkje korleis, og det gjev ingen måte å avgjere om ein utøvar reflekterer godt eller dårleg. Schön sjølv visste dette og gjorde det til ein dygd: kunnskapen er taus, seier han, han lèt seg ikkje fullt artikulere. Men eit fag som gjer det \emph{uartikulerbare} til kjernen sin, har erklært at kjernen ikkje kan delast — og udelbar kunnskap er per definisjon ikkje-kumulativ. Schön gav faget ein generøs orsaking for ikkje å måtte seie kva det veit. At verket er sitert på tvers av medisin, jus, pedagogikk og design, viser styrken i orsakinga, ikkje styrken i fundamentet.

Richard Buchanan sine «wicked problems», sitert kring åtte tusen gonger, er den andre \autocite{buchanan1992wicked}. Vi har vist korleis ein presis observasjon — at nokre problem manglar endeleg form — vart vridd til ein identitet: at designens \emph{vesen} er det vrange, og at faget difor ikkje \emph{kan} ha eit kumulativt fundament. Her skal berre poenget om sitering leggjast til. Kvifor vart nett denne artikkelen så uhyrleg sitert? Ikkje fordi han løyste noko, men fordi han gav faget den mest verdfulle gåva av alle: ei prinsipiell grunngjeving for å sleppe å produsere det andre fag må produsere. Du kan ikkje klandre eit fag for å mangle kumulative resultat når faget har lært, frå sitt mest siterte verk, at problema dets per natur ikkje gjev slike. Artikkelen er sitert fordi han avvæpnar kritikken. Det er den fullkomne sjølvstadfestande teksten: jo meir faget vert utfordra på fråvêret av eit fundament, jo flittigare siterer det forklaringa på kvifor fråvêret er nødvendig.

Klaus Krippendorff sin «semantic turn» er den tredje, og han er den mest ambisiøse av alle, for han ber undertittelen «ein ny grunnvoll for design» \autocite{krippendorff2006,krippendorff1984semantics}. Det er nett kva boka leitar etter — eit eksplisitt forsøk på å reise fundamentet — og difor er dommen tung. Krippendorff hevda at design framføre alt er å gje ting meining, og at faget burde organiserast kring meiningsskaping. Den sterkaste forma er reell: gjenstandar tyder verkeleg noko, og funksjonalismen fortrengte det. Men ein grunnvoll må gje ein målestokk, og meininga fekk ingen. Når tydinga ligg hjå alle som omgår gjenstanden, kan han bety alt for nokon, og då har faget framleis ingen måte å seie at éi tyding er rettare enn ei anna. Det er det mest avslørande av alt at mannen som meir enn nokon hadde levert ein \emph{prøvbar} metode i nabofaget — innhaldsanalysen, med sin reliabilitetskoeffisient — gav design ein «grunnvoll» heilt utan ein slik. Krippendorff visste presis kva ein etterprøvbar målestokk var. Han bygde ein i samfunnsvitskapen. Og då han skreiv om design, lét han vere å byggje ein, ikkje av vankunne, men fordi materialet ikkje bar det. Skilnaden mellom hans to bøker er skilnaden mellom eit fag som har eit fundament og eit som berre har eit vokabular.

Tre verk, tre grunnsteinar, og det same i kvar: ein verkeleg innsikt som endar i noko uprøvbart — det tause, det vrange, det grenselaust tydbare. Faget peikar på desse tre når det vil vise at det har eit fundament. Men les ein dei som det dei er, peikar dei alle same vegen: mot grunnen til at fundamentet aldri kom.

\section{Historikarane, sosialinnovatørane og lånet utanfrå}

To grupper står att, og dei rundar av mønsteret heller enn å bryte det.

Den eine er dei som tok faget si manglande grunngjeving som eit politisk og historisk emne heller enn eit filosofisk. Victor Margolin, sitert kring åtte tusen gonger, gav design ei historieskriving og ein plass i det politiske — han grunnla \emph{Design Issues}, skreiv ein verdshistorie om design, kravde at faget såg seg sjølv som samfunnskraft \autocite{margolin2002politics,margolin2015,margolin1992histories}. Ezio Manzini, lengre nede i tala men kanoniserte for det, flytta design mot sosial innovasjon og berekraft \autocite{manzini2015}. Harold Nelson og Erik Stolterman gav faget «the design way», eit forsøk på å heve design til ein tredje kultur ved sida av vitskap og humaniora \autocite{nelson2012designway,stolterman2008nature}. Gjev dei den sterkaste forma: dette er seriøst, lærd, ofte vakkert arbeid, og verda er betre for at nokon skreiv det. Men sjå kva det \emph{ikkje} gjer. Margolin gjev faget ei historie og ein etikk; han gjev det ikkje ein teori om form, og han hevda heller ikkje å gjere det. Manzini gjev faget eit oppdrag; eit oppdrag er ikkje eit fundament, det er ein retning å sakne eitt i. Nelson og Stolterman gjev faget eit \emph{krav} på å vere ein eigen kultur — men eit krav er ikkje ei oppfylling. Heile denne gruppa er innflytelsesrik fordi ho gjev faget meining, retning og sjølvkjensle. Ingen av delane er det faget manglar. Det faget manglar, er ein prøvbar kjerne, og om den skriv denne gruppa, klokt nok, ikkje.

Den andre er lånet i sin reinaste form: Christopher Frayling. Hans vesle skrift frå 1993 gav faget skiljet mellom forsking \emph{inn i}, \emph{for} og \emph{gjennom} design — den mest siterte ordninga i heile diskusjonen om kva designforsking er \autocite{frayling1993}. Og her finst det eit funn så talande at det kunne stått som motto over heile kapitlet. Granskarar av Frayling-resepsjonen har peika på at langt over nitti prosent av dei som siterer, omtalar eller parafraserer skiljet hans, aldri har lese originalen — eit kort, vanskeleg tilgjengeleg skrift som dei færraste har hatt i handa. Faget siterer altså si eiga mest grunnleggjande formulering om forsking på andrehand, frå parafrase av parafrase, utan kjeldekontakt. Står det att eit klårare bilete av kva sitering tyder i dette faget? Eit prøvbart resultat treng ikkje lesast på tru; det kan etterprøvast. Ei innflytelsesrik formulering treng ikkje lesast i det heile; ho treng berre vidareformidlast, til ho er overalt og ingen lenger hugsar kva ho eigentleg sa. Frayling-sitatet er ikkje overføring av kunnskap. Det er forplantinga av eit namn.

Samlar ein dei ti — Norman, Krippendorff, Simon på lån utanfrå; Cross, Dorst, Lawson for kunnskapsforma; Schön, Buchanan for orsakingane; Margolin og Manzini, Nelson og Stolterman, Frayling for historia, oppdraget og forskingsretorikken — ser ein at ikkje ein einaste av dei vart mest sitert for å ha lagt ein prøvbar stein i ein mur andre kunne byggje vidare på. Dei vart siterte for å ha sagt noko treffande, frigjerande eller forklarande om kva design \emph{er}. Det er ekte innflytelse. Men innflytelse av denne sorten akkumulerer ikkje; han berre spreier seg. Ein kan vere den mest siterte i eit fag og likevel ikkje ha gjeve det eit fundament — og når \emph{alle} dei mest siterte er det, er kanon ikkje eit motbevis mot fråvêret. Kanon \emph{er} fråvêret, ført i pennen av dei flinkaste i rommet.

\section{Kvifor toppen stadfestar tomrommet}

Det står att å seie kvifor dette ikkje berre er ein samling uheldige einskildtilfelle, men ein nødvendig følgje av kva slag fag design er. Lat oss vere heilt rettferdige mot namna éin siste gong: ingen av dei feila fordi dei var dumme. Dei var, dei fleste, blant dei skarpaste hovuda faget har hatt, og fleire av dei var skarpare enn nesten alle i nabofaga. Det er nett det som gjer saka tvingande. Når dei beste i eit fag, over femti år, med all den intelligensen og lærdomen feltet kunne mønstre, samla ikkje produserer ein einaste prøvbar, kumulativ kjerne — men i staden produserer den eine vakre, uprøvbare formuleringa etter den andre — då er ikkje forklaringa at dei var utilstrekkelege. Forklaringa er at materialet ikkje bar det dei prøvde å reise. Eit fag der toppen er full av geni og tom for fundament, fortel deg at fundamentet ikkje var mogleg å nå med dei midlane faget hadde.

Her ligg sjølve diagnosen, og ho er hardare enn ho fyrst ser ut. I eit fag med eit fundament er dei mest siterte verka resultat: Watson og Crick sin struktur, som ti tusen seinare arbeid byggjer på og foredlar; ein teori som vart prøvd, haldt og utvida. Sitatet peikar bakover til noko som vart \emph{vist}, og framover til noko som vert \emph{bygd}. I designfaget peikar dei mest siterte verka ingen av delane. «Designerly ways of knowing», «the reflective practitioner», «wicked problems», «the semantic turn» — dei vart ikkje viste, og dei vart ikkje bygde vidare på i den forstand at noko vart \emph{tilført} og prøvd. Dei vart sagde, og så sagde att, og så att, til dei var sjølvsagte. Det er forskjellen mellom ein mur og eit ekko. Ein mur veks når kvar stein ber den neste. Eit ekko veks ikkje; det berre held fram med å lyde, høgare di fleire vegger det kastast mellom, til alle høyrer det og ingen kan seie kvar det kom frå.

Difor er det mest siterte i designfaget det mest avslørande, ikkje det minst. Innvendinga vi opna med — at store namn må tyde eit fundament — har det presis bakvendt. Eit fag med eit fundament treng ikkje dyrke namna sine, fordi det har resultata deira; det kan gløyme kven som fyrst sa ein ting, fordi tingen står att, prøvd og innbygd, uavhengig av munnen. Eit fag utan fundament har berre munnane. Det må hugse kven som sa, fordi det \emph{er} ingenting utan kven som sa det — inga prøve som overlever utan opphavsmannen, ingen stein som ber seg sjølv. Designfaget kanoniserer fordi det ikkje kan akkumulere. Det held fast på dei ti namna med ein iver som forrår at det ikkje har noko anna å halde fast på. Og det er, til sjuande og sist, kvifor dei er ti namn og ikkje ti resultat: eit resultat treng ikkje eit namn for å bestå, men ei meining gjer det — ho døyr i det augneblinket ingen lenger gidd sitere henne.

\vignett

Men om sjølv toppen av faget berre stadfestar fråvêret, då melder eit siste spørsmål seg, og det er ikkje lenger eit spørsmål om fortida: kva med eit fag som veit alt dette om seg sjølv, og likevel held fram med å dele ut gradar i det? Det er sjølvbedraget som no skal granskast — ikkje som villfaring, men som ein stabil og innretta likevekt.

\vidarelesing{\fullcite{cross2006}; \fullcite{schon1983}; \fullcite{buchanan1992wicked}; \fullcite{krippendorff2006}; \fullcite{dorst2011core}; \fullcite{frayling1993}; \fullcite{nelson2012designway}.}
